% !TeX root = ../main.tex

\chapter{绪论}
\section{研究背景与意义}
近年来高考数学愈发强调试题的综合性, 重视知识间的联系, 也更注重学生数学思维的考察, 而含参函数的“任意性”与“存在性”问题正符合这一特点, 一些研究把它称之为高中函数方程与不等式的综合问题, 更常见的叫法是函数的“恒成立”与“能成立”问题, 这类问题主要考察学生对量词和函数的理解, 通过全称量词和存在量词的约束将函数方程与不等式问题转化为函数求最值或值域的问题, 从而限制并求解参数的范围。这类问题可以从函数个数、变量个数和量词选择三个方面来进行分类, 渗透着数形结合、化归与转化、函数与方程等许多重要数学思想, 解法也层出不穷, 如判别式法、数形结合法、导数分析法、分离参数法等, 具有很强的逻辑性, 对提高学生的逻辑思维能力和综合应用能力有很大的帮助, 研究此问题具有一定的创新价值与理论意义。

从教学大纲和考试要求来看, 含参函数的“任意性”与“存在性”问题在各类考试中都有一定的占比, 而在高考数学中, 这类问题常常出现在函数综合题、 不等式与函数结合的题目中, 如 24 年新课标 II 卷第 8 题、23 年全国甲卷 (理) 的第 21 题的第(2)小题、23 年新课标 I 卷的第 19 题的第(2)小题等, 是高考命题的热点。

然而, 由于其逻辑复杂性和综合性, 这类问题成为了教学中的重点和难点内容。教师在教学过程中往往发现学生常因概念模糊、逻辑混乱而难以掌握这类问题, 而且这些学生占比非常大, 这也导致研究它的教学意义相对变小, 这也是相关研究比较少的原因。传统教学多采用题型加套路的方式, 虽能应付考试, 但不利于学生理解本质, 难以灵活运用。因此我们需要深入研究题目本质及其考察重点, 研究其解题方法思路并优化教学策略, 将复杂问题简单化, 技巧性问题通俗化, 帮助学生更好地理解此类问题, 提高教学质量, 满足教学需求。

目前关于含参函数的“任意性”与“存在性”问题的研究大多集中在双函数双变量等单一题型的研究, 缺乏系统地针对函数个数、变量个数和量词选择三个方面分类的题型的全面研究。大多数研究也仅限于总结部分题型和解题策略, 缺少针对此类问题的教学研究。我们需要深入探讨如何从函数本质出发, 让学生理解量词、变量和函数之间的逻辑关系, 构建完善的教学体系, 因此研究此问题具有一定的创新价值与实践意义。


\section{研究现状}
当前关于含参函数“任意性”与“存在性”问题的研究主要集中在双函数双变量等单一题型的解题方法探讨方面, 缺少全面的题型梳理, 以及系统的教学策略的研究。

在理论研究方面，学者们普遍关注此类问题的解题技巧。徐国平 \cite{2007_徐国平_用函数最值法处理“恒成立”不等式中的参数范围}、郑良 \cite{2011_郑良_反思解题过程变通解题方法——有感于分离参数法与函数最值法的对话}、 赖鸿杰 \cite{2012_赖鸿杰_最值法在高中数学恒成立问题中的应用} 等指出最值法是重要解题方法，通过求函数最值或值域，将不等式与函数相结合，还可通过分类讨论处理不确定因素；王利华 \cite{2012_王利华_分离变量是基本通法}、刘元德 \cite{2019_刘元德_含参不等式问题的三种求解策略}、谭芳芳 \cite{2021_谭芳芳_“分参法”妙解导数解答题} 强调分离变量和参数分离策略能简化问题；严波研究了利用判别式法解决一元二次不等关系恒成立问题 \cite{2014_严波_不等式恒成立问题的求解策略}；刘元德 \cite{2019_刘元德_含参不等式问题的三种求解策略}、王雪芹 \cite{2014_王雪芹_恒成立问题及存在性问题的解题策略}、厉立军 \cite{2019_厉立军_转换主元解放思想} 提出分类讨论和转换主元的思路；杨苍洲阐述了数形结合思想及其应用局限 \cite{2020_杨苍洲_函数含参问题的速解策略——参数半分离}；成银玉介绍了构造函数法，包括直接构造和间接构造 \cite{2014_成银玉_构造函数、利用导数解答不等式恒成立问题}。这些方法相互补充，形成了一套完整的解题体系，能够解决各类恒成立问题。然而这些研究多侧重于单一题型的解法分析，缺乏对不同题型间内在联系的系统性研究。

在教育实践方面，诸多学者提出了各具特色的教学思路。陆斌借助创设学校雕塑情境及设计变式问题，激发学生兴趣，提升思维品质\cite{2024_陆斌_“5E学习环”模式下的高中数学复习教学——以一类函数恒成立问题教学为例}；何智强调“教学评一体化”, 通过明确目标、设计问题串及即时评价来促进教学 \cite{2022_何智_“教学评一体化”视域下的高中数学课堂教学探索与实践——以高三复习课“不等式恒成立问题”为例}；李德安和孙雪梅的提出“套课”教学模式，分阶段提升学生思维 \cite{2015_李德安_高中数学解题教学“套课”教学模式的探讨}；付高南 \cite{2024_付高南_高中数学解题教学中学生高阶思维能力的培养——以恒成立求参数问题为例}、于素娟 \cite{2024_于素娟_高中数学解题教学中学生高阶思维能力的培养——以恒成立求参数问题为例} 主张通过开放式解题、融入数学思想和变式训练培养学生高阶思维；江战明注重剖析“成立”概念，避免学生解题出错 \cite{2013_江战明_何为成立何为恒成立——一个不等式证明引发的教学思考}；侯斐斐、冯炜基于核心素养开发的“问题阶梯”教学法有效提升了教学效果 \cite{2022_侯斐斐_基于核心素养的探究课教学——以求解函数中的双变量“任意性”和“存在性”问题为例}；李刚基于单元教学观，整合内容、渗透思想并评价教学效果 \cite{2022_李刚_基于单元教学观的高三数学复习案例设计——以“不等式恒成立问题”为例}。这些研究为含参函数“任意性”与“存在性”问题的教学提供了丰富的思路和方法借鉴，然而这些研究是只针对单一题型的教学，没有涉及到所有题型的系统化的教学。


教材编写方面，以人教版和北师大版为代表的国内教材采用循序渐进的编排方式。在高一至高二上的函数与导数的基础学习阶段，人教 A 版的必修一“二次函数”章节中，通过根的分布问题初步渗透这类问题，此时一般用判别式法和分类讨论法求解；在人教 A 版的选择性必修二“导数应用”中，涉及到用最值法求此类问题，做法是将此类问题转化为求函数的最值或值域问题，此时利用导数工具求函数最值或值域即可求解此类问题。在高二下的专题深化阶段出现“恒成立”专题，此时涉及到一些题型及方法总结，强调解题方法的系统归纳。在高三的综合应用复习阶段会涉及到实际背景的复杂应用题，综合性较强，强调与高考的衔接。

近年来随着新课程改革的推进，诸多研究聚焦于教学方法与技术工具对学生学习效果的影响，其中 GeoGebra 软件的应用备受关注。陈守月、多布杰提出数形结合可以优化教师教学并促进学生学习，而 Geogebra 是实现数形结合的强有力的工具 \cite{2025_陈守月_GeoGebra在数学教育中的创新应用与价值研究}。Ndagijimana 等人通过对卢旺达学生的研究发现，在社会文化学习情境中运用 GeoGebra 软件，能有效提升学生对数学概念的理解，特别是在点、线、角等知识的学习上，实验组学生成绩显著优于对照组\cite{2024_Ndagijimana_Contributions_of_GeoGebra_software_within_the_socio-cultural_proximity_on_enhancing_students’conceptual_understanding_of_mathematics}。而 Romero Albaladejo 和 García López 在西班牙开展的教学实验表明, GeoGebra 软件有助于改善学生的数学态度，如增强学生的毅力、提高精确严谨性以及培养自主性\cite{2024_Romero_Albaladejo_Mathematical_attitudes_transformation_when_introducing_GeoGebra_in_the_secondary_classroom}。 这些研究揭示了 GeoGebra 软件在数学教育中的积极作用，为数学教育工作者提供了极具价值的参考，也为后续相关研究奠定了坚实基础。笔者想利用 GeoGebra 软件将数形结合融入此类专题教学，帮助学生理解题意。然而，现有研究在这类问题的现代教育技术与传统教学的融合创新等方面仍存在明显不足，这为后续研究指明了发展方向。


\section{研究思路与方法}

本研究采用理论分析与教学设计相结合的研究思路, 将研究内容明确划分为题型系统分类和教学方案设计两个部分。

在题型分类研究中, 这类问题的分类方法丰富多样, 大体分为按解法和题目条件两种分类方法, 按解法分类的有判别式法、分类讨论法、数形结合法、导数分析法、分离参数法和最值法等类别, 本文按照题目条件分类, 具体按照所给函数个数、变量个数和量词的组合形式交叉分类, 大体分为单函数单变量、单函数双变量、双函数单变量和双函数双变量四大类, 在这四大类下又细分为任意型、 存在型, 以及任意-任意型、任意-存在型和存在-存在型。

每种题型下会给出题型的一般形式和一般解法思路, 以及一道典型例题, 所有题型所给的例题都由一道题目变化而来, 此题目为笔者实习时高一学生问的练习题, 原题具体为本文给出的双函数双变量中任意-存在型的例题。此分类方法以核心函数为出发点, 通过系统调整条件生成各类题型, 建立了统一的函数模型变式体系, 仅保证了知识体系的完整覆盖, 同时也明确了不同题型之间的本质联系, 相对完整地分析了含参函数的“任意性”与“存在性”问题。

在完成题型分类后, 本研究又设计了新的教学方案。具体做法是: 用一个工厂营收的实际问题作为案例引入, 设计一个连接所有题型教学的问题链, 同样建立统一的函数模型变式体系, 帮助学生形成系统性知识; 把问题镶嵌在实际问题背景中, 让学生把现实问题转化成数学问题, 帮助学生提高解决实际问题的能力。 然后使用 GeoGebra 软件动态展示参数变化对函数图像的影响, 让学生深入理解题目变化规律, 引导学生将题目进行转化, 训练学生数形结合、转化与化归等数学思想。

本文教学设计大体分为单函数和双函数类型的教学两部分, 由单函数问题引入专题, 双函数问题深入专题。对于单函数问题, 由老师示范解题过程, 双函数问题则让学生分组合作, 通过 GeoGebra 软件调整参数、观察图像变化, 自己探索解题方法。

这种“教师引导加自主探究”的设计, 既保证了概念传递的准确性, 又通过小组合作培养了学生的自主学习和深入思考能力。借助技术工具, 抽象的数学问题变得直观易懂, 帮助学生真正掌握知识, 并能灵活运用到实际问题中。